% ==============================================================================
% PG - Lorenzo Rizzi Fiorot
% Capítulo 2 - Referencial Teórico e Tecnologias Utilizadas
% ==============================================================================
\chapter{Referencial Teórico e Tecnologias Utilizadas}
\label{sec-referencial}

% TODO(autor): no anteprojeto, a abertura do capítulo também apresentava o repositório, como
% relação entre a BNCC e os pilares do Pensamento Computacional; essa apresentação depende
% do enquadramento do trabalho.
Este capítulo reúne os fundamentos teóricos do trabalho e as tecnologias empregadas em seu desenvolvimento. A Seção~\ref{sec-referencial-bncc} apresenta a estrutura da BNCC e da BNCC Computação, e a Seção~\ref{sec-referencial-pc}, o Pensamento Computacional. A Seção~\ref{sec-referencial-msl} caracteriza o mapeamento sistemático da literatura como método de pesquisa. A Seção~\ref{sec-referencial-delimitacao} delimita o conceito de repositório de conhecimento adotado. A Seção~\ref{sec-referencial-visualizacao} trata da visualização de grafos, e a Seção~\ref{sec-referencial-acessibilidade}, da acessibilidade na web. Por fim, a Seção~\ref{sec-referencial-tecnologias} descreve as tecnologias utilizadas.


%%% Início de seção. %%%
\section{A BNCC e a BNCC Computação}
\label{sec-referencial-bncc}

A BNCC é um ``documento de caráter normativo que define o conjunto orgânico e progressivo de aprendizagens essenciais que todos os alunos devem desenvolver ao longo das etapas e modalidades da Educação Básica'' \cite[p.~7]{brasil2018bncc}. O documento abrange a Educação Infantil, o Ensino Fundamental e o Ensino Médio e constitui referência nacional para a formulação dos currículos dos sistemas e das redes de ensino e das propostas pedagógicas das escolas \cite[p.~8--9]{brasil2018bncc}. As aprendizagens essenciais devem concorrer para o desenvolvimento de dez competências gerais, entendida a competência como ``a mobilização de conhecimentos (conceitos e procedimentos), habilidades (práticas, cognitivas e socioemocionais), atitudes e valores para resolver demandas complexas da vida cotidiana, do pleno exercício da cidadania e do mundo do trabalho'' \cite[p.~8]{brasil2018bncc}.

No Ensino Fundamental, etapa a que pertencem os dados do grafo apresentado neste trabalho, a BNCC organiza-se em cinco áreas do conhecimento --- Linguagens, Matemática, Ciências da Natureza, Ciências Humanas e Ensino Religioso ---, que reúnem os componentes curriculares \cite[p.~27]{brasil2018bncc}. Cada componente apresenta um conjunto de habilidades, que expressam as aprendizagens essenciais a serem asseguradas aos alunos. As habilidades relacionam-se a objetos de conhecimento, isto é, a conteúdos, conceitos e processos, e estes, por sua vez, organizam-se em unidades temáticas \cite[p.~28--29]{brasil2018bncc}. Uma unidade temática contempla uma gama maior ou menor de objetos de conhecimento, e cada objeto relaciona-se a um número variável de habilidades \cite[p.~29]{brasil2018bncc}. Esse encadeamento --- área do conhecimento, componente curricular, unidade temática, objeto de conhecimento e habilidade --- é a estrutura reproduzida no modelo do grafo.

Cada habilidade é identificada por um código alfanumérico \cite[p.~30]{brasil2018bncc}, cuja composição é apresentada na Tabela~\ref{tab:codigo-habilidade}, tomando-se como exemplo o código \texttt{EF01CI05}. Duas ressalvas do documento são relevantes para este trabalho: a numeração sequencial não representa uma ordem ou hierarquia esperada das aprendizagens, e a organização das habilidades em objetos de conhecimento e unidades temáticas expressa ``um arranjo possível (dentre outros)'', que não deve ser tomado como modelo obrigatório para o desenho dos currículos \cite[p.~31]{brasil2018bncc}. Ambas se estendem ao grafo, que reproduz esse arranjo.

\begin{table}[h]
    \centering
    \caption{Composição do código das habilidades do Ensino Fundamental}
    \label{tab:codigo-habilidade}
    \renewcommand{\arraystretch}{1.3}
    \begin{tabular}{|p{4.2cm}|p{6cm}|p{4.2cm}|}
        \hline
        \textbf{Parte do código} & \textbf{O que indica} & \textbf{No código \texttt{EF01CI05}} \\
        \hline
        Primeiro par de letras & Etapa de ensino & \texttt{EF}: Ensino Fundamental \\
        \hline
        Primeiro par de números & Ano ou bloco de anos & \texttt{01}: 1º ano \\
        \hline
        Segundo par de letras & Componente curricular & \texttt{CI}: Ciências \\
        \hline
        Último par de números & Posição na numeração sequencial do ano ou do bloco de anos & \texttt{05}: quinta habilidade \\
        \hline
    \end{tabular}
\end{table}

A Computação já era tratada na BNCC de 2018, que tematiza três dimensões da computação e das tecnologias digitais \cite[p.~473--474]{brasil2018bncc}: o pensamento computacional, discutido na Seção~\ref{sec-referencial-pc}; o mundo digital, que envolve as formas de processar, transmitir e distribuir a informação de maneira segura e confiável; e a cultura digital, voltada a uma participação mais consciente e democrática por meio das tecnologias digitais. No Ensino Fundamental, contudo, essas dimensões foram contempladas nas competências e habilidades dos diferentes componentes \cite[p.~474]{brasil2018bncc}, sem habilidades próprias, uma vez que a Computação não figura entre os componentes curriculares da etapa \cite[p.~27]{brasil2018bncc}.

O conjunto próprio de habilidades veio com a Resolução CNE/CEB nº 1, de 4 de outubro de 2022, que define normas sobre Computação na Educação Básica, em complemento à BNCC (art.~1º). A Resolução determina que o desenvolvimento e a formulação dos currículos considerem as tabelas de competências e habilidades anexas (art.~1º, §~2º), atribui aos Estados, aos Municípios e ao Distrito Federal a definição dos parâmetros e das abordagens pedagógicas de implementação (art.~2º) e prevê que o Ministério da Educação (MEC), em conjunto com esses entes, defina política para a formação docente, o apoio ao desenvolvimento de currículos e o apoio ao desenvolvimento de recursos didáticos compatíveis com as tabelas (art.~4º) \cite{brasil2022computacao}.

As tabelas constam do documento anexo ao Parecer CNE/CEB nº 2/2022 \cite{brasil2022complemento} e cobrem as três etapas da Educação Básica. No Ensino Fundamental, há sete competências e um conjunto de habilidades para cada ano, além de habilidades por etapa (do 1º ao 5º e do 6º ao 9º ano). Cada habilidade vincula-se a um objeto de conhecimento e a um dos três eixos --- Pensamento Computacional, Mundo Digital e Cultura Digital --- e vem acompanhada de explicação e de exemplos. Os códigos seguem a composição da Tabela~\ref{tab:codigo-habilidade}, com o par \texttt{CO} na posição do componente curricular: \texttt{EF01CO01}, por exemplo, identifica a primeira habilidade de Computação do 1º ano, ``Organizar objetos físicos ou digitais considerando diferentes características para esta organização, explicitando semelhanças (padrões) e diferenças'' \cite[p.~12]{brasil2022complemento}. A organização das tabelas é, portanto, análoga à dos demais componentes, com o eixo no lugar da unidade temática.

% TODO(autor): se o enquadramento do trabalho incluir a BNCC Computação no modelo do grafo,
% acrescentar aqui a relação entre as tabelas do complemento e o modelo.


%%% Início de seção. %%%
\section{Pensamento Computacional}
\label{sec-referencial-pc}

A popularização do termo Pensamento Computacional é atribuída ao artigo de \citeonline{wing2006}, como registram \citeonline[p.~562]{boeres2024}. Para a autora, trata-se de uma habilidade fundamental para todos, e não apenas para cientistas da computação, a ser acrescentada à leitura, à escrita e à aritmética na formação de toda criança; envolve a resolução de problemas, o projeto de sistemas e a compreensão do comportamento humano com base nos conceitos fundamentais da ciência da computação \cite[p.~33]{wing2006}. Nessa acepção, pensar computacionalmente significa conceituar, e não programar, e constitui um modo pelo qual as pessoas, e não os computadores, resolvem problemas \cite[p.~35]{wing2006}.

No contexto brasileiro, \citeonline{brackmann2017} propõe, a partir da fusão de diversas fontes, a seguinte definição:

\begin{citacao}
O Pensamento Computacional é uma distinta capacidade criativa, crítica e estratégica humana de saber utilizar os fundamentos da Computação, nas mais diversas áreas do conhecimento, com a finalidade de identificar e resolver problemas, de maneira individual ou colaborativa, através de passos claros, de tal forma que uma pessoa ou uma máquina possam executá-los eficazmente \cite[p.~29]{brackmann2017}.
\end{citacao}

\noindent A BNCC, por sua vez, afirma que o pensamento computacional ``envolve as capacidades de compreender, analisar, definir, modelar, resolver, comparar e automatizar problemas e suas soluções, de forma metódica e sistemática, por meio do desenvolvimento de algoritmos'' \cite[p.~474]{brasil2018bncc}. As três formulações tratam o PC como um modo de resolver problemas fundamentado na Computação, e as de \citeonline{wing2006} e \citeonline{brackmann2017} ressaltam que os processos envolvidos podem ser executados por uma pessoa ou por uma máquina. Por essa razão, o PC pode ser desenvolvido também por meio da Computação Desplugada, que dispensa o uso de computadores e recorre, por exemplo, a movimentos, cartões, desenhos e enigmas \cite[p.~50]{brackmann2017}.

Segundo \citeonline[p.~33]{brackmann2017}, o PC apoia-se em quatro pilares, interdependentes e voltados à resolução de problemas:

\begin{itemize}
    \item \textbf{decomposição}: identificar um problema complexo e dividi-lo em partes menores e mais fáceis de gerenciar;
    \item \textbf{reconhecimento de padrões}: analisar cada uma dessas partes em busca de problemas parecidos, já solucionados anteriormente;
    \item \textbf{abstração}: concentrar-se apenas nos detalhes importantes, ignorando as informações irrelevantes;
    \item \textbf{algoritmos}: criar passos ou regras simples para resolver cada um dos subproblemas encontrados.
\end{itemize}

\noindent Os mesmos pilares são adotados por \citeonline[p.~565]{boeres2024}, na análise da formação docente, e por \citeonline[p.~1752]{osorio2025}, no mapeamento entre o PC e as habilidades da BNCC.

Nos documentos curriculares, o PC aparece de duas formas. Na BNCC de 2018, as menções no Ensino Fundamental concentram-se na área de Matemática: os processos de resolução de problemas, de investigação, de desenvolvimento de projetos e de modelagem são considerados potencialmente ricos para o seu desenvolvimento \cite[p.~266]{brasil2018bncc}, e ressaltam-se a importância dos algoritmos e de seus fluxogramas e a identificação de padrões para estabelecer generalizações \cite[p.~271]{brasil2018bncc}. Na BNCC Computação, o PC constitui um eixo próprio, e os pilares manifestam-se nas habilidades desde os anos iniciais: no 1º ano, a identificação de semelhanças, ou padrões, na organização de objetos (\texttt{EF01CO01}) e a criação de sequências de passos, associadas à palavra ``Algoritmos'' (\texttt{EF01CO03}); no 2º ano, a identificação de atributos essenciais, própria da abstração, na modelagem de objetos (\texttt{EF02CO01}); no 3º ano, a decomposição de problemas complexos (\texttt{EF03CO03}) \cite[p.~12, 16, 18]{brasil2022complemento}.

% TODO(autor): parágrafo de fechamento, que liga o Pensamento Computacional ao repositório
% de conhecimento; depende do enquadramento do trabalho.


%%% Início de seção. %%%
\section{Mapeamento Sistemático da Literatura}
\label{sec-referencial-msl}

O Mapeamento Sistemático da Literatura (MSL) é uma metodologia de pesquisa secundária que visa identificar, categorizar e sintetizar o conjunto de evidências disponíveis sobre um determinado tema, distinguindo-se da revisão sistemática convencional por priorizar uma visão panorâmica da área em vez de responder a uma questão clínica ou empírica específica \cite{petersen2008systematic}. Enquanto a revisão sistemática tende a focar em perguntas de pesquisa bem delimitadas e a realizar sínteses quantitativas dos resultados, o mapeamento sistemático permite classificar e estruturar o estado da arte de forma mais abrangente, evidenciando lacunas, tendências e distribuição de estudos ao longo do tempo \cite{kitchenham2007guidelines, petersen2008systematic}. No presente trabalho, essa abordagem foi adotada para levantar o estado da arte sobre o ensino do Pensamento Computacional no contexto da educação básica brasileira, com ênfase na integração com os componentes curriculares da BNCC, fornecendo a base empírica necessária para a construção do repositório de conhecimento proposto. O protocolo, a condução e os resultados do mapeamento são apresentados no Capítulo~\ref{sec-msl}.


%%% Início de seção. %%%
\section{Delimitação Conceitual: Repositório de Conhecimento}
\label{sec-referencial-delimitacao}

A literatura técnica e científica emprega diferentes designações para soluções baseadas na web, cujas distinções conceituais são relevantes para a correta caracterização do artefato desenvolvido neste trabalho. Por \textit{repositório} entende-se, no campo da Ciência da Informação e da Arquivística, uma solução voltada ao armazenamento, ao gerenciamento e à preservação de longo prazo de materiais digitais, constituída por hardware, software, metadados e procedimentos normativos \cite{conarq2015}. Na Engenharia de Software, o termo designa ainda o local em que um sistema de controle de versão registra e mantém o histórico de alterações do código-fonte, como nos repositórios Git \cite{chacon2014}. Em ambas as acepções, a ênfase recai sobre a guarda e a organização de um acervo --- de conteúdo ou de código ---, e não sobre a sua mera apresentação.

Distinguem-se desse conceito outras designações correntes. Um \textit{site} (ou \textit{website}) é um conjunto de páginas web interligadas e hospedadas em um mesmo endereço de rede, acessadas a partir de uma página inicial \cite{w3c1999}, noção centrada na disponibilização de conteúdo e com interatividade tipicamente limitada. Um \textit{portal} designa um ponto único de acesso que agrega, em uma mesma interface, informações, recursos e serviços de fontes diversas, funcionando como porta de entrada para a utilização de serviços na internet \cite{gouveia2007}; sua característica definidora é a agregação, e não a estruturação de um acervo de conhecimento. Já uma \textit{plataforma}, em sua acepção mais ampla, corresponde a uma infraestrutura digital que possibilita a interação entre dois ou mais grupos de usuários \cite{srnicek2017}, conceito associado a ecossistemas extensíveis e de natureza sociotécnica, cuja aplicação exige rigor de definição \cite{dereuver2018}. O termo \textit{ambiente}, por sua vez, remete a um espaço de execução ou de interação e, no domínio educacional, consolidou-se na expressão \textit{ambiente virtual de aprendizagem} (AVA), que designa espaços digitais de mediação do processo de ensino-aprendizagem apoiados em tecnologias de informação e comunicação \cite{almeida2003}.

Cabe, por fim, distinguir os termos \textit{sistema web} e \textit{aplicação web} (\textit{webapp}). O primeiro designa um software hospedado em servidor e acessado por navegador, geralmente restrito a usuários cadastrados, em contraste com o site meramente informativo \cite{sommerville2011}. O segundo, mais preciso para soluções que integram interface de usuário (\textit{front-end}) e lógica de servidor (\textit{back-end}), compreende desde sites completos até funcionalidades especializadas e aplicações de processamento de dados \cite{pressman2011}.

Diante dessas distinções, este trabalho adota o termo \textit{repositório de conhecimento} para designar o artefato central produzido, uma vez que nenhuma das demais designações (\textit{site}, \textit{portal}, \textit{plataforma} ou \textit{ambiente}) captura adequadamente a natureza de acervo organizado e semanticamente explícito da solução. Esse repositório materializa-se na forma de um grafo de conhecimento. O termo \textit{aplicação web} é empregado, de modo complementar e não substitutivo, para designar especificamente a camada de software que disponibiliza o repositório, reservando-se o termo \textit{repositório} ao acervo de conhecimento assim disponibilizado.


%%% Início de seção. %%%
\section{Visualização de Grafos}
\label{sec-referencial-visualizacao}

Texto.


\subsection{Grafos e o Seu Desenho}
\label{sec-referencial-visualizacao-desenho}

Texto.

\begin{figure}[H]
    \centering
    \caption{O mesmo recorte do grafo desenhado pelo layout por forças e pelo layout hierárquico}
    \label{fig-vis-layouts}
    \captura{figuras/fig-vis-layouts.jpg}
    \notafigura{print já feito: o mesmo grafo (1º ano, ciências, ``o ano inteiro nesta matéria'') no layout de forças, na esquerda, e no hierárquico, na direita. recortei só o desenho, sem o painel e sem a legenda.}
\end{figure}


\subsection{Layout por Forças}
\label{sec-referencial-visualizacao-forcas}

Texto.


\subsection{Layout Hierárquico}
\label{sec-referencial-visualizacao-hierarquico}

Texto.

\begin{figure}[H]
    \centering
    \caption{Etapas do layout hierárquico}
    \label{fig-vis-etapas-hierarquico}
    \marcadorbrasao
    \notafigura{aqui vai um esquema simples, feito no próprio latex, com as etapas do hierárquico: separar os nós em camadas, ordenar cada camada pra cruzar menos arestas e depois calcular a posição de cada nó. podem ser três caixinhas com setas.}
\end{figure}



%%% Início de seção. %%%
\section{Acessibilidade na Web}
\label{sec-referencial-acessibilidade}

Texto.


\subsection{Diretrizes de Acessibilidade para Conteúdo Web}
\label{sec-referencial-acessibilidade-wcag}

Texto.

\begin{table}[H]
    \centering
    \caption{Critérios das WCAG atendidos pelo EduGraph}
    \label{tab:wcag}
    \renewcommand{\arraystretch}{1.3}
    \begin{tabular}{|l|l|l|}
        \hline
        \textbf{Critério} & \textbf{O que pede} & \textbf{Como o site atende} \\
        \hline
        Texto. & Texto. & Texto. \\
        \hline
    \end{tabular}
    \notafigura{aqui vai uma tabela com os critérios das wcag que o site atende: contraste, não depender só da cor, aumentar o texto, usar pelo teclado, foco visível e animação que dá pra desligar. cada linha com o número do critério, o que ele pede e onde isso aparece no site.}
\end{table}


\subsection{Deficiências na Percepção de Cores}
\label{sec-referencial-acessibilidade-cores}

Texto.

\begin{figure}[H]
    \centering
    \caption{Cores dos níveis do grafo na visão típica e nas simulações de daltonismo}
    \label{fig-acess-simulacao}
    \marcadorbrasao
    \notafigura{aqui vai uma tabela de cores, feita no latex mesmo, com a cor de cada nível do grafo (ano, área, componente, unidade, objeto e habilidade) na visão normal e simulando protanopia, deuteranopia e tritanopia. a ideia é mostrar que dá pra diferenciar as cores em todas.}
\end{figure}



%%% Início de seção. %%%
\section{Tecnologias Utilizadas}
\label{sec-referencial-tecnologias}

Esta seção apresenta as tecnologias adotadas no desenvolvimento do EduGraph e as razões de cada escolha; a Tabela~\ref{tab:tecnologias} resume-as, com as versões utilizadas, e a forma como se articulam é descrita na Seção~\ref{sec-edugraph-arquitetura}.

No \textit{front-end}, adotou-se o React, biblioteca JavaScript de código aberto voltada à construção de interfaces de usuário a partir de componentes declarativos e reutilizáveis. O modelo de componentes ajusta-se à decomposição natural da interface projetada, a barra de busca e filtros, o palco de visualização do grafo, o painel de detalhe das entidades e as sobreposições de estado, permitindo que cada unidade encapsule estrutura, aparência e comportamento próprios, o que favorece o desenvolvimento incremental e a manutenção isolada de cada parte. A renderização declarativa complementa esse modelo: a interface é descrita como função do estado da aplicação, de modo que transições como consulta em andamento, resultado disponível, resultado vazio e falha da consulta refletem-se automaticamente na tela, sem manipulação manual da árvore do documento. Pesou ainda o caráter de biblioteca e não de \textit{framework} do React, que não impõe arquitetura de aplicação, solução de roteamento ou camada de estilos, preservando as demais decisões do projeto.

Em torno do React organizam-se as demais escolhas da camada de apresentação. O Vite atua como empacotador, respondendo pelo servidor de desenvolvimento e pela geração dos artefatos de produção. A estilização emprega CSS (\textit{Cascading Style Sheets}) puro, organizado em variáveis de tema que centralizam a paleta nos modos claro e escuro de exibição, sem o recurso a \textit{frameworks} de componentes prontos --- o que assegura controle integral sobre a identidade visual construída na etapa de prototipação, descrita na Seção~\ref{sec-edugraph-prototipacao}, e sobre os recursos de acessibilidade da ferramenta. Os estilos combinam declarações em linha nos componentes, para a estrutura e para as propriedades calculadas em função do estado da aplicação, e uma folha central com as variáveis de tema e as classes responsáveis pelos estados de interação e pelas animações. A navegação entre as telas também é controlada pelo estado do React, sem biblioteca de roteamento, de modo que as únicas dependências de execução da interface são os pacotes \texttt{react} e \texttt{react-dom}. A visualização do grafo, por sua vez, é desenhada em um elemento \textit{canvas} do HTML (\textit{HyperText Markup Language}), por meio de sua interface de desenho bidimensional (Canvas 2D), com um motor próprio de simulação e de interação, sem bibliotecas de visualização de grafos, descrito na Seção~\ref{sec-edugraph-grafos}.

A integração entre as camadas reforça essas escolhas. A interface consulta a API do próprio sistema por meio de requisições HTTP (\textit{Hypertext Transfer Protocol}), com respostas em JSON (\textit{JavaScript Object Notation}), recursos nativos do navegador utilizados a partir dos componentes React, sem bibliotecas de acesso específicas. O acoplamento entre a interface e o servidor restringe-se, assim, ao contrato das consultas e das respostas, o que permitiu, nas fases iniciais do desenvolvimento, substituir o servidor por um simulador que reproduzia o contrato então previsto, viabilizando a evolução da interface antes da disponibilização dos dados.

No armazenamento, adotou-se o Neo4j, banco de dados em grafo cuja linguagem declarativa de consulta é a Cypher \cite{angles2017foundations}. O acesso ao banco se dá pela Query API, interface que executa instruções Cypher recebidas em requisições HTTP e devolve o resultado em JSON \cite{neo4j2026queryapi}. A opção por essa interface, em detrimento dos \textit{drivers} oficiais do banco, permitiu adotar em todo o sistema um único formato de resposta, o JSON tipado, no qual cada valor é acompanhado da indicação do próprio tipo, como nó, relação, caminho, inteiro, entre outros, que a API repassa sem conversão e a partir do qual a interface reconstrói o grafo de qualquer resultado de consulta.

Entre a interface e o banco, adotou-se uma API própria, escrita em Python com o \textit{framework} FastAPI e executada pelo servidor Uvicorn; a comunicação com o Neo4j usa a biblioteca httpx, e a configuração de acesso ao banco é lida de um arquivo local pela biblioteca python-dotenv. A escolha do FastAPI considerou a validação dos dados recebidos por meio de tipos, a geração automática da documentação interativa das rotas, que favorece a continuidade do projeto por outros pesquisadores, e o suporte a rotas assíncronas, adequado a serviços que respondam de forma incremental. A mesma API deverá receber os serviços previstos para o sistema --- o fórum, o repositório de ontologias e o assistente baseado em modelos de linguagem ---, cujos módulos-base já constam dela (Seção~\ref{sec-edugraph-arquitetura}).

\begin{table}[h]
    \centering
    \caption{Tecnologias utilizadas e respectivas versões}
    \label{tab:tecnologias}
    \renewcommand{\arraystretch}{1.3}
    \begin{tabular}{|>{\raggedright\arraybackslash}p{4.2cm}|>{\raggedright\arraybackslash}p{2.4cm}|>{\raggedright\arraybackslash}p{7.6cm}|}
        \hline
        \textbf{Tecnologia} & \textbf{Versão} & \textbf{Uso no projeto} \\
        \hline
        React (\texttt{react} e \texttt{react-dom}) & 18.3.1 & Construção da interface web \\
        \hline
        Vite & 5.4.21 & Servidor de desenvolvimento e geração da versão de produção da interface \\
        \hline
        Node.js & 22.22.1 & Execução das ferramentas de desenvolvimento da interface \\
        \hline
        Python & 3.14.4 & Linguagem da API \\
        \hline
        FastAPI & 0.165.0 & \textit{Framework} da API: rotas, validação dos dados e documentação \\
        \hline
        Uvicorn & 0.54.0 & Servidor que executa a API \\
        \hline
        httpx & 0.28.1 & Cliente HTTP da API para o Neo4j \\
        \hline
        python-dotenv & 1.2.4 & Leitura da configuração de acesso ao banco \\
        \hline
        Neo4j, edição Community & 2026.09.0 & Banco de dados em grafo \\
        \hline
    \end{tabular}
\end{table}
